\chapter{Trabajo relacionado}
\label{sec:paper-related}

GraphHunter se ubica en la intersección de tres líneas: lenguajes
de consulta de SIEM, \emph{threat hunting} sobre grafos de
procedencia, y detección de anomalías sobre grafos heterogéneos.
Este capítulo discute representativos de cada línea y sitúa las
decisiones de diseño propias respecto a ellos.

\section{Lenguajes de consulta en SIEM}

Los SIEM contemporáneos ofrecen lenguajes orientados a tablas con
extensiones para series de tiempo: Splunk SPL, Azure KQL, Elastic
EQL. La expresividad es suficiente para \emph{detection rules} y
búsquedas textuales con agregaciones. La estructura relacional se
modela con joins textuales sucesivos, legibles en ejemplos
pequeños, frágiles en investigaciones de varios hops.

La diferencia conceptual es \emph{dónde vive} la estructura del
incidente. En SIEM, vive en la redacción de la query (`... | join
... | join ...`) que el analista debe mantener coherente. En
GraphHunter, vive en el grafo canonical y la hipótesis DSL
(cap.~\ref{sec:paper-dfs}) sólo describe el patrón; la
materialización de paths la hace el motor.

\section{Threat hunting sobre grafos de procedencia}

HOLMES~\cite{milajerdi2019holmes} y Poirot~\cite{milajerdi2019poirot}
inauguraron la línea de construir grafos causales a partir de
eventos del sistema operativo (fork/exec, open, connect) y
emparejarlos contra patrones derivados de ATT\&CK. RapSheet y
trabajos posteriores refinan la forma de los patrones y las
heurísticas de ranking. El target común es el grafo de procedencia
del host.

GraphHunter difiere en tres ejes respecto a esta línea:

\begin{enumerate}
    \item \textbf{Fuente de datos}. En lugar de grafo causal del
    kernel, trabaja sobre un grafo canonical construido desde logs
    heterogéneos (OCSF v1.4). Permite razonar entre hosts y entre
    dominios (red + IAM + cloud audit), no sólo dentro de un host.
    \item \textbf{Lenguaje}. El DSL es externo y
    composicional~(cap.~\ref{sec:paper-dfs}); HOLMES/Poirot
    embeben los patrones en reglas internas. Permite a un analista
    escribir hipótesis nuevas sin tocar código.
    \item \textbf{Scoring}. Componemos cinco señales de anomalía
    operables (cap.~\ref{sec:paper-score}). La línea HOLMES tiende
    a scoring binario (match/no-match) con umbrales ocultos.
\end{enumerate}

\section{Detección de anomalías sobre grafos heterogéneos}

Una línea paralela emplea representaciones aprendidas
(DeepWalk~\cite{perozzi2014deepwalk}, embeddings GNN) para
detectar anomalías sin patrón explícito. Las arquitecturas con GCN
sobre grafos de logs son efectivas cuando existe volumen de datos
etiquetados. Tienen dos costos operativos bien documentados:
distribución shift entre entornos (un modelo entrenado en un banco
no generaliza a una telco) y dificultad de dar explicaciones al
analista.

GraphHunter usa anomalía como \emph{señal composicional}: los cinco
componentes (entity rarity, edge rarity, neighborhood
concentration, temporal novelty, gnn threat) son
interpretables por separado. El componente GNN está tras un
\emph{model gate} (cap.~\ref{sec:paper-gate}) y no bloquea el uso
del sistema sin modelo entrenado. Esto es deliberado: el sistema
debe funcionar el primer día en un cliente sin dataset propio.

\section{Superficies agénticas}

La interacción con asistentes LLM via Model Context Protocol (MCP)
es un espacio todavía en estandarización. Los trabajos publicados
se enfocan en servidores MCP genéricos sobre bases de datos
tabulares o APIs REST. GraphHunter expone 55 tools agrupadas en 10
dominios semánticos (search, graph, hunt, catalog,
explorer, tags, notes, export, ingest, sessions) permitiendo a
Claude Code o Cursor conducir investigaciones completas.

\section{Compiladores asistidos por LLM}

El ingestor agéntico (cap.~\ref{sec:paper-ingestor}) explota una
idea conocida en \emph{program synthesis}: el LLM propone un
artefacto de \emph{configuración}, no el código ejecutable. El
código lo produce un compilador determinístico. Esa separación
preserva la propiedad de determinismo productivo que el LLM por sí
sólo no da. El patrón es similar al usado en
Copilot-for-SQL y variantes: el LLM sugiere, un validador rechaza
lo inválido, el humano ratifica.
